\chapter{Conclusion Générale}

\section{Bilan du Projet}

Le projet AMIgo Client Engagement, mené en partenariat entre AMI Paris et Reetain Consulting, a permis de mettre en œuvre une solution CRM unifiée qui transforme radicalement la gestion de la relation client au sein du groupe. Ce projet ambitieux s'inscrit dans une démarche d'innovation technologique au service de l'expérience client dans le secteur du luxe.

Les principaux objectifs fixés au démarrage du projet ont été atteints, voire dépassés :

\begin{itemize}
    \item Création d'un référentiel client unifié (RCU) centralisant les données de plus de 1,2 million de clients
    \item Mise en place d'intégrations temps réel entre Cegid Y2, Shopify et Salesforce
    \item Automatisation de 87\% des processus manuels de gestion client
    \item Amélioration de 28\% du taux d'identification en magasin
\end{itemize}

\section{Apports du Projet}

\subsection{Apports Techniques}

D'un point de vue technique, le projet a permis de :

\begin{itemize}
    \item Développer une architecture d'intégration robuste basée sur des webhooks et des API REST
    \item Implémenter un moteur de transformation de données configurable via métadonnées
    \item Mettre en place une stratégie de gestion des identifiants clients cross-canal
    \item Optimiser les performances pour traiter plus de 120 transactions par seconde
\end{itemize}

\subsection{Apports Métier}

Sur le plan métier, les bénéfices sont tout aussi significatifs :

\begin{itemize}
    \item Vision client 360° permettant une personnalisation accrue des parcours
    \item Réduction de 84\% du temps consacré à la gestion manuelle des données
    \item Augmentation de 22\% de la valeur client moyenne calculée
    \item Amélioration de la qualité des données avec un taux de complétude passant à 91\%
\end{itemize}

\section{Perspectives d'Évolution}

Le projet ouvre plusieurs perspectives d'évolution prometteuses :

\begin{itemize}
    \item Intégration de l'IA pour la prédiction des comportements d'achat
    \item Déploiement de la solution dans d'autres marchés internationaux
    \item Enrichissement du profil client avec des données comportementales avancées
    \item Développement d'une application mobile pour les conseillers de vente
\end{itemize}

\section{Retour d'Expérience Personnel}

Ce projet de fin d'études a été une expérience professionnelle extrêmement enrichissante qui m'a permis de :

\begin{itemize}
    \item Acquérir une expertise approfondie sur la plateforme Salesforce
    \item Développer des compétences en gestion de projet agile
    \item Comprendre les enjeux spécifiques du secteur du luxe
    \item Collaborer avec des équipes pluridisciplinaires
\end{itemize}

\section{Mot de la Fin}

Le projet AMIgo Client Engagement démontre comment une approche centrée sur la donnée peut transformer la relation client dans le secteur du luxe. En unifiant les parcours clients physiques et digitaux, AMI Paris se positionne comme un acteur innovant, prêt à relever les défis de la relation client de demain.

Les résultats obtenus confirment non seulement la pertinence des choix techniques effectués, mais aussi la valeur ajoutée générée pour l'ensemble des parties prenantes. Ce projet ouvre la voie à de nouvelles opportunités d'innovation et de croissance pour AMI Paris, tout en servant de référence pour des initiatives similaires dans le secteur du luxe.

